\subsection{Cost of \appname}
\begin{table}[t]
\centering
\caption{Average cost of \appname and baselines across models.}
\label{tab:cost_bound}
\footnotesize
\setlength{\tabcolsep}{5pt}
\renewcommand{\arraystretch}{0.95}
\resizebox{\columnwidth}{!}{
\begin{tabular}{lcccc}
\toprule
Method 
& Loc. time (s) 
& Edit time (s) 
& Test time (s) 
& Artifact size \\
\midrule
\cellcolor{boundRow}\textbf{\appname}
& \cellcolor{boundRow}23.81 
& \cellcolor{boundRow}83.09 
& \cellcolor{boundRow}11.30 
& \cellcolor{boundRow}731.03 KB \\

ROME
& 10.19 
& 384.45 
& 10.86 
& 98.01 MB \\

MEMIT
& -- 
& 193.06 
& 13.17 
& 588.01 MB \\

DINM
& 0.92 
& 11.97 
& 12.63 
& 126.68 MB \\

Full-FT
& -- 
& 124.16 
& 15.31 
& 14.27 GB \\

\scriptsize{Self-Refinement}
& -- 
& --  
& 25.12 
& --  \\
\bottomrule
\end{tabular}
}
\vspace{-0.3cm}
\end{table}

% To assess the practical cost of \appname, we evaluate its runtime and storage overhead. Table~\ref{tab:cost_bound} reports the average results across the three evaluated models.
% \appname introduces a localization step that takes 23.81 seconds on average. Although this adds extra cost, RQ3 shows that localization helps identify modules related to the package-validity boundary and improves editing effectiveness.
% For editing, \appname requires 83.09 seconds on average, which is lower than ROME (384.45s), MEMIT (193.06s), and Full-FT (124.16s).
% At inference time, \appname takes 11.30 seconds per case on average, which is comparable to ROME (10.86s) and lower than MEMIT, DINM, Full-FT, and Self-Refinement.
% Its average artifact size is only 731.03 KB, much smaller than ROME (98.01 MB), MEMIT (588.01 MB), DINM (126.68 MB), and Full-FT (14.27 GB).
% This is because \appname stores only LoRA updates on localized modules rather than large parameter edits or full checkpoints.
% Overall, \appname adds a small localization cost but achieves efficient editing, competitive inference time, and very small storage overhead. 
% These results show that \appname is practical for mitigating package hallucination.

To assess the practical cost of \appname, we evaluate its runtime and storage overhead. 
Table~\ref{tab:cost_bound} reports the average results across the three evaluated models. 
\appname introduces a localization step that takes 23.81 seconds on average. 
Although this adds extra cost, RQ3 shows that localization improves editing effectiveness by identifying modules related to the package-validity boundary.
For editing, \appname requires 83.09 seconds on average, which is lower than ROME, MEMIT, and Full-FT. 
At inference time, \appname takes 11.30 seconds per case on average, which is comparable to ROME and lower than other baselines. 
\appname is also lightweight in storage, requiring only 731.03 KB of artifacts on average, approximately 20,000× smaller than full fine-tuning checkpoints. This is because \appname stores only LoRA updates on localized modules rather than large parameter edits or full checkpoints. 
Overall, \appname adds a small localization cost while achieving efficient editing, competitive inference time, and very small storage overhead.
These results show that \appname is practical for mitigating package hallucination.

\subsection{Impact of Editing on Model Performance}
\begin{table}[t]
\centering
\captionsetup{skip=2pt}
\caption{Code accuracy on HumanEval before and after editing.}
\label{tab:editing_drop}
\footnotesize
\setlength{\tabcolsep}{3pt}
\renewcommand{\arraystretch}{0.95}
\resizebox{\columnwidth}{!}{
\begin{tabular}{lccc|ccc}
\toprule
\multirow{2}{*}{Model}
& \multicolumn{3}{c|}{pass@1}
& \multicolumn{3}{c}{pass@10} \\
\cmidrule(lr){2-4} \cmidrule(lr){5-7}
& Orig. & Edited & $\Delta$
& Orig. & Edited & $\Delta$ \\
\midrule
DeepSeekCoder
& 68.96\% & 67.58\% & -1.38\%
& 82.32\% & 83.39\% & +1.07\% \\

Qwen3
& 23.29\% & 21.93\% & -1.36\%
& 51.83\% & 50.91\% & -0.92\% \\

Llama-3.1
& 41.52\% & 42.12\% & +0.60\%
& 69.51\% & 69.66\% & +0.15\% \\
\bottomrule
\end{tabular}
}
\vspace{-0.3cm}
\end{table}

While \appname effectively reduces package hallucination, it is important to examine whether editing affects the model's general code generation ability. 
To this end, we evaluate the original and edited models on HumanEval. 
Following the main evaluation setup, each model has four edited versions, obtained by editing on the four edit folds separately. 
We evaluate all four edited models on HumanEval and report the average pass@1 and pass@10 results. 
Table~\ref{tab:editing_drop} presents the results.

The results show that \appname remains stable on clean code generation inputs, with most changes within a $\pm$2\% margin. 
Overall, \appname mitigates package hallucination with minimal impact on the original capabilities of the model.


\subsection{Threats to Validity}
\noindent\textbf{Internal Validity:}
A potential internal threat comes from package extraction. 
We rely on regular expressions to extract package names from model outputs. 
Due to the inherent instability of LLM outputs, the extracted packages may not always follow the expected format. 
To reduce this threat, we constrain the output format in the prompt and use one-shot prompting to guide generation.
Another concern is the stability of evaluation. 
Since LLM generation involves randomness, the measured hallucination rates may exhibit slight variance. 
We address this by generating five responses for each prompt, and reporting average results over four edit folds.
The module localization step of \appname may also introduce uncertainty. 
Although \appname identifies modules related to package boundaries using a risk-aware localization objective, the results may still be affected by noisy gradients or the selected edit cases. 
To mitigate this threat, we aggregate localization signals over multiple edit cases. 

\noindent\textbf{External Validity:}
Our evaluation is conducted on three open-source LLMs from different model families. However, the results may not directly generalize to other LLMs. We focus on open-source models because they provide stable access after release, allowing historical versions to be revisited and reproduced. In contrast, commercial LLMs are continuously updated and may become unavailable over time.
Another external threat comes from the package ecosystem studied in this paper. 
We focus on Python packages from PyPI because PyPI provides public metadata for package validity checking. 
Future work can extend \appname to additional package ecosystems and programming languages (e.g., Java) to further evaluate its generality.
Furthermore, we evaluate \appname on three package-related tasks: package recommendation, code generation, and \texttt{pip install} recommendation. Although these tasks cover common package-related development scenarios, real-world dependency management also involves configuration files, build scripts, and repository-level workflows. Future work can extend \appname to these broader settings.